% GENERATED FILE. Edit canonical lessons, then run npm run generate:books.
% canonical-source-hash: fnv1a64:31dfd94532fc4994
% canonical-lessons: ML-C528-sathyam, ML-C528-kathinadhvani, ML-C528-subhaptivisvasam, ML-C528-atmartthata, ML-C528-accatakkam

\chapter{Stubborn and Hard-working}
\label{ch:ml-stubborn-and-hard-working}
\hlchaptermodality{528}

\begin{chapteropening}
\textbf{By the end of this chapter:} \emph{I can say stubbornness, a hard worker, optimism, sincerity and discipline.}

Complete the last lesson of chapter 528: I can say stubbornness, a hard worker, optimism, sincerity and discipline.
\end{chapteropening}

\section[\texorpdfstring{\ml{ശാഠ്യം} (śāṭhyaṁ)}{ശാഠ്യം (śāṭhyaṁ)}]{\ml{ശാഠ്യം} (śāṭhyaṁ) — stubbornness}

\label{lesson:ML-C528-sathyam}



\begin{quote}
Before the new one: say the Malayalam for generous, then the Malayalam for a loyal person.
\end{quote}

\subsection*{You'll want to know: \ml{ശാഠ്യം}}
\textbf{\ml{ശാഠ്യം}} — \emph{śāṭhyaṁ} — \textquotedblleft{}stubbornness\textquotedblright{}.

\begin{quote}
\ml{കളിപ്പാട്ടം} \ml{വേണമെന്ന്} \ml{കുട്ടി} \ml{ശാഠ്യം} \ml{പിടിച്ചു}.
\end{quote}

\begin{quote}
\emph{kaḷippāṭṭaṁ vēṇamennŭ kuṭṭi śāṭhyaṁ piṭiccu.}
\end{quote}

\begin{quote}
The child stubbornly insisted on having a toy.
\end{quote}

\ml{ശാഠ്യം} \ml{പിടിക്കുക}, to dig in one's heels, is the usual phrase.

\begin{cousinweb}
\textbf{\ml{ശാഠ്യം}} (\emph{śāṭhyaṁ}) is borrowed from Sanskrit \emph{śāṭhya}, \textquotedblleft{}wickedness, guile\textquotedblright{} (Monier-Williams). In Malayalam the word came to mean a stubborn refusal to give way.
\end{cousinweb}

\begin{grammarlens}[title={What you've built}]
One more word for describing what you lived and felt.
\end{grammarlens}

\subsection*{Your turn}
\begin{itemize}
\raggedright
  \item \emph{Say it:} \emph{śāṭhyaṁ}
  \item \emph{Say it:} \emph{śāṭhyaṁ}, once more
  \item \emph{Say it:} \emph{kaḷippāṭṭaṁ vēṇamennŭ kuṭṭi śāṭhyaṁ piṭiccu.} — The child stubbornly insisted on having a toy.
  \item \emph{Recall:} say the Malayalam for generous, then the Malayalam for a loyal person, then say \emph{śāṭhyaṁ} again
\end{itemize}

\begin{tcolorbox}[breakable,colback=teal!4,colframe=teal!35!black,title={Before you move on}]
What does \ml{ശാഠ്യം} mean? (\textquotedblleft{}Stubbornness\textquotedblright{}.) Say it once more.
\end{tcolorbox}
\section[\texorpdfstring{\ml{കഠിനാധ്വാനി} (kaṭhinādhvāni)}{കഠിനാധ്വാനി (kaṭhinādhvāni)}]{\ml{കഠിനാധ്വാനി} (kaṭhinādhvāni) — a hard worker}

\label{lesson:ML-C528-kathinadhvani}



\begin{quote}
Before the new one: say the Malayalam for a loyal person, then the Malayalam for stubbornness.
\end{quote}

\subsection*{You'll want to know: \ml{കഠിനാധ്വാനി}}
\textbf{\ml{കഠിനാധ്വാനി}} — \emph{kaṭhinādhvāni} — \textquotedblleft{}a hard worker\textquotedblright{}.

\begin{quote}
\ml{എന്റെ} \ml{അമ്മ} \ml{ഒരു} \ml{കഠിനാധ്വാനിയാണ്}.
\end{quote}

\begin{quote}
\emph{enṟe amma oru kaṭhinādhvāniyāṇŭ.}
\end{quote}

\begin{quote}
My mother is a hard worker.
\end{quote}

\ml{കഠിനം} hard + \ml{അധ്വാനം} labour + -\ml{ഇ}, a person. Hard work itself is \ml{കഠിനാധ്വാനം}.

\begin{cousinweb}
\textbf{\ml{കഠിനാധ്വാനി}} (\emph{kaṭhinādhvāni}) is a Malayalam compound built from Sanskrit \emph{kaṭhina}, \textquotedblleft{}hard\textquotedblright{} (CDIAL 2650), and \ml{അധ്വാനം}, \textquotedblleft{}labour\textquotedblright{}, with the ending -\ml{ഇ} for a person. Where \ml{അധ്വാനം} itself comes from is not settled.
\end{cousinweb}

\begin{grammarlens}[title={What you've built}]
One more word for describing what you lived and felt.
\end{grammarlens}

\subsection*{Your turn}
\begin{itemize}
\raggedright
  \item \emph{Say it:} \emph{kaṭhinādhvāni}
  \item \emph{Say it:} \emph{kaṭhinādhvāni}, once more
  \item \emph{Say it:} \emph{enṟe amma oru kaṭhinādhvāniyāṇŭ.} — My mother is a hard worker.
  \item \emph{Recall:} say the Malayalam for a loyal person, then the Malayalam for stubbornness, then say \emph{kaṭhinādhvāni} again
\end{itemize}

\begin{tcolorbox}[breakable,colback=teal!4,colframe=teal!35!black,title={Before you move on}]
What does \ml{കഠിനാധ്വാനി} mean? (\textquotedblleft{}A hard worker\textquotedblright{}.) Say it once more.
\end{tcolorbox}
\section[\texorpdfstring{\ml{ശുഭാപ്തിവിശ്വാസം} (śubhāptiviśvāsaṁ)}{ശുഭാപ്തിവിശ്വാസം (śubhāptiviśvāsaṁ)}]{\ml{ശുഭാപ്തിവിശ്വാസം} (śubhāptiviśvāsaṁ) — optimism}

\label{lesson:ML-C528-subhaptivisvasam}



\begin{quote}
Before the new one: say the Malayalam for stubbornness, then the Malayalam for a hard worker.
\end{quote}

\subsection*{You'll want to know: \ml{ശുഭാപ്തിവിശ്വാസം}}
\textbf{\ml{ശുഭാപ്തിവിശ്വാസം}} — \emph{śubhāptiviśvāsaṁ} — \textquotedblleft{}optimism\textquotedblright{}.

\begin{quote}
\ml{എല്ലാം} \ml{ശരിയാകുമെന്ന്} \ml{അവന്} \ml{ശുഭാപ്തിവിശ്വാസമുണ്ട്}.
\end{quote}

\begin{quote}
\emph{ellāṁ śariyākumennŭ avanŭ śubhāptiviśvāsamuṇṭŭ.}
\end{quote}

\begin{quote}
He is optimistic that everything will be all right.
\end{quote}

\ml{ശുഭം} good + \ml{ആപ്തി} reaching + \ml{വിശ്വാസം} trust: trust that things will turn out well.

\begin{cousinweb}
\textbf{\ml{ശുഭാപ്തിവിശ്വാസം}} (\emph{śubhāptiviśvāsaṁ}) is a learned compound built from Sanskrit \emph{śubha}, \textquotedblleft{}good, auspicious\textquotedblright{} (CDIAL 12532), \emph{āpti}, \textquotedblleft{}reaching, gain\textquotedblright{} (Monier-Williams), and \emph{viśvāsa}, \textquotedblleft{}trust\textquotedblright{} (CDIAL 11966).
\end{cousinweb}

\begin{grammarlens}[title={What you've built}]
One more word for describing what you lived and felt.
\end{grammarlens}

\subsection*{Your turn}
\begin{itemize}
\raggedright
  \item \emph{Say it:} \emph{śubhāptiviśvāsaṁ}
  \item \emph{Say it:} \emph{śubhāptiviśvāsaṁ}, once more
  \item \emph{Say it:} \emph{ellāṁ śariyākumennŭ avanŭ śubhāptiviśvāsamuṇṭŭ.} — He is optimistic that everything will be all right.
  \item \emph{Recall:} say the Malayalam for stubbornness, then the Malayalam for a hard worker, then say \emph{śubhāptiviśvāsaṁ} again
\end{itemize}

\begin{tcolorbox}[breakable,colback=teal!4,colframe=teal!35!black,title={Before you move on}]
What does \ml{ശുഭാപ്തിവിശ്വാസം} mean? (\textquotedblleft{}Optimism\textquotedblright{}.) Say it once more.
\end{tcolorbox}
\section[\texorpdfstring{\ml{ആത്മാർത്ഥത} (ātmārtthata)}{ആത്മാർത്ഥത (ātmārtthata)}]{\ml{ആത്മാർത്ഥത} (ātmārtthata) — sincerity}

\label{lesson:ML-C528-atmartthata}



\begin{quote}
Before the new one: say the Malayalam for a hard worker, then the Malayalam for optimism.
\end{quote}

\subsection*{You'll want to know: \ml{ആത്മാർത്ഥത}}
\textbf{\ml{ആത്മാർത്ഥത}} — \emph{ātmārtthata} — \textquotedblleft{}sincerity\textquotedblright{}.

\begin{quote}
\ml{അവളുടെ} \ml{ജോലിയിലെ} \ml{ആത്മാർത്ഥത} \ml{എല്ലാവർക്കും} \ml{ഇഷ്ടമാണ്}.
\end{quote}

\begin{quote}
\emph{avaḷuṭe jōliyile ātmārtthata ellāvarkkuṁ iṣṭamāṇŭ.}
\end{quote}

\begin{quote}
Everyone likes the sincerity of her work.
\end{quote}

The adjective is \ml{ആത്മാർത്ഥമായ}, sincere.

\begin{cousinweb}
\textbf{\ml{ആത്മാർത്ഥത}} (\emph{ātmārtthata}) is built from Sanskrit \emph{ātman}, \textquotedblleft{}self, soul\textquotedblright{}, and \emph{artha}, \textquotedblleft{}aim, purpose\textquotedblright{} (Monier-Williams), with the noun ending -\ml{ത}: acting from one's own heart.
\end{cousinweb}

\begin{grammarlens}[title={What you've built}]
One more word for describing what you lived and felt.
\end{grammarlens}

\subsection*{Your turn}
\begin{itemize}
\raggedright
  \item \emph{Say it:} \emph{ātmārtthata}
  \item \emph{Say it:} \emph{ātmārtthata}, once more
  \item \emph{Say it:} \emph{avaḷuṭe jōliyile ātmārtthata ellāvarkkuṁ iṣṭamāṇŭ.} — Everyone likes the sincerity of her work.
  \item \emph{Recall:} say the Malayalam for a hard worker, then the Malayalam for optimism, then say \emph{ātmārtthata} again
\end{itemize}

\begin{tcolorbox}[breakable,colback=teal!4,colframe=teal!35!black,title={Before you move on}]
What does \ml{ആത്മാർത്ഥത} mean? (\textquotedblleft{}Sincerity\textquotedblright{}.) Say it once more.
\end{tcolorbox}
\section[\texorpdfstring{\ml{അച്ചടക്കം} (accaṭakkaṁ)}{അച്ചടക്കം (accaṭakkaṁ)}]{\ml{അച്ചടക്കം} (accaṭakkaṁ) — discipline}

\label{lesson:ML-C528-accatakkam}



\begin{quote}
Before the new one: say the Malayalam for optimism, then the Malayalam for sincerity.
\end{quote}

\subsection*{You'll want to know: \ml{അച്ചടക്കം}}
\textbf{\ml{അച്ചടക്കം}} — \emph{accaṭakkaṁ} — \textquotedblleft{}discipline\textquotedblright{}.

\begin{quote}
\ml{ആ} \ml{സ്കൂളിൽ} \ml{കർശനമായ} \ml{അച്ചടക്കം} \ml{ഉണ്ടായിരുന്നു}.
\end{quote}

\begin{quote}
\emph{ā skūḷil karśanamāya accaṭakkaṁ uṇṭāyirunnu.}
\end{quote}

\begin{quote}
That school had strict discipline.
\end{quote}

Its second part is \ml{അടക്കം}, restraint. \ml{അച്ചടക്കമുള്ള} means disciplined.

\begin{cousinweb}
\textbf{\ml{അച്ചടക്കം}} (\emph{accaṭakkaṁ}) is a native compound whose second part, \ml{അടക്കം}, \textquotedblleft{}holding down, restraint\textquotedblright{}, is clear (DEDR 63); the origin of the first part is not settled.
\end{cousinweb}

\begin{grammarlens}[title={What you've built}]
One more word for describing what you lived and felt.
\end{grammarlens}

\subsection*{Your turn}
\begin{itemize}
\raggedright
  \item \emph{Say it:} \emph{accaṭakkaṁ}
  \item \emph{Say it:} \emph{accaṭakkaṁ}, once more
  \item \emph{Say it:} \emph{ā skūḷil karśanamāya accaṭakkaṁ uṇṭāyirunnu.} — That school had strict discipline.
  \item \emph{Recall:} say the Malayalam for optimism, then the Malayalam for sincerity, then say \emph{accaṭakkaṁ} again
\end{itemize}

\begin{tcolorbox}[breakable,colback=teal!4,colframe=teal!35!black,title={Before you move on}]
What does \ml{അച്ചടക്കം} mean? (\textquotedblleft{}Discipline\textquotedblright{}.) Say it once more.
\end{tcolorbox}
